% Modul Belajar 7, dihasilkan oleh tools/build.py dari tools/konten/p07.py
\documentclass[a4paper]{article}
\usepackage{modulITERA}
\begin{document}
\Sampul{7}{Review dan Simulasi UTS}{Merangkai materi Pertemuan 1 sampai 6 dan berlatih dengan soal bergaya UTS}{CPMK0607 (menerapkan) dan CPMK0608 (menjelaskan) pada konsep dasar algoritma dan sistem komputer}{sekitar 3 jam belajar mandiri, ditambah 150 menit di kelas}{\texttt{02 Hands-on/P7 - Review UTS - Hands-on/}}
\Bagian{Modul 7}{Peta belajar}
Ikuti urutan ini dari atas ke bawah. Setiap bagian memakai apa yang sudah dibahas di bagian sebelumnya.
\par\vspace{4pt}\noindent{\fontsize{9.5pt}{13pt}\selectfont
\begin{tabular}{@{}>{\raggedright\arraybackslash}p{0.140\textwidth}>{\raggedright\arraybackslash}p{0.840\textwidth}@{}}
\kepala{Urutan} & \kepala{Bagian}\\
1 & Tentang modul ini\\\garis
2 & Masalah pembuka: Satu soal, semua materi\\\garis
3 & Langkah 1: Peta materi UTS\\\garis
4 & Langkah 2: Strategi soal A\\\garis
5 & Langkah 3: Strategi soal B\\\garis
6 & Langkah 4: Contoh tracing gabungan\\\garis
7 & Langkah 5: Memecah digit\\\garis
8 & Pesan error yang sering muncul\\\garis
9 & Latihan mandiri\\\garis
10 & Rangkuman\\\garis
11 & Cek pemahaman\\\garis
12 & Exit Ticket\\\garis
13 & Bacaan lanjut dan persiapan minggu depan\\
\garisdasar
\end{tabular}}\par\vspace{6pt}
\Bagian{Pembuka}{Tentang modul ini}
Modul ini mendampingi Pertemuan 7. Tidak ada materi baru. Tujuannya merangkai semua yang sudah dipelajari dan berlatih dengan soal bergaya UTS: soal A menulis program, soal B menelusuri dan menjelaskan.

Isinya sama dengan slide di kelas, tetapi ditulis lebih pelan dan lebih lengkap, supaya kalian bisa mengulang sendiri di rumah tanpa harus menebak maksud slide.

\Sub{Setelah menyelesaikan modul ini, kalian bisa}
\par\vspace{4pt}\noindent{\fontsize{9.5pt}{13pt}\selectfont
\begin{tabular}{@{}>{\raggedright\arraybackslash}p{0.070\textwidth}>{\raggedright\arraybackslash}p{0.680\textwidth}>{\raggedright\arraybackslash}p{0.230\textwidth}@{}}
\kepala{No} & \kepala{Kemampuan} & \kepala{Dibahas di}\\
1 & Menyelesaikan masalah yang menggabungkan masukan, ekspresi, percabangan, dan perulangan dalam satu program C++ dalam batas waktu ujian (Sub-CPMK 7.1, CPMK0607) & Langkah 1 sampai 5\\\garis
2 & Menelusuri program yang menggabungkan percabangan dan perulangan bersarang, lalu menjelaskan setiap perubahan nilai dengan istilah yang tepat (Sub-CPMK 7.2, CPMK0608) & kotak Tebak dulu, Latihan 3\\
\garisdasar
\end{tabular}}\par\vspace{6pt}
Karena setiap instrumen penilaian dibagi rata antara Bagian A (menulis kode) dan Bagian B (menelusuri dan menjelaskan), yang dinilai bukan hanya program kalian jalan, tetapi juga kalian bisa menjelaskan mengapa program itu jalan.

\Sub{Cara memakai modul ini}
Setiap langkah punya pola yang sama: penjelasan konsep, kadang contoh dari kehidupan sehari-hari, lalu kode yang bisa langsung kalian ketik. Sesudah kode ada kotak \textbf{Tebak dulu}. Tulis tebakan kalian di kertas sebelum membaca \textbf{Jawaban} di bawahnya. Kebiasaan menebak lalu membuktikan inilah yang membuat konsep menempel, bukan sekadar membaca.

\begin{Penting}[Ketik, jangan salin]
Ketik ulang setiap contoh di GDB Online atau editor kalian, jangan disalin-tempel. Saat mengetik, kalian akan membuat salah ketik, membaca pesan error, dan memperbaikinya. Itu bagian dari belajar.
\end{Penting}
\Sub{Yang perlu disiapkan}
Peramban dengan akses ke onlinegdb.com, atau g++ di komputer kalian sendiri. Kertas dan alat tulis untuk menebak keluaran dan membuat trace table. Folder hands-on pertemuan ini dari repo mata kuliah.

\Bagian{Pembuka}{Masalah pembuka: Satu soal, semua materi}
Soal UTS jarang menguji satu topik saja. Contohnya tarif parkir: membaca masukan, membulatkan menit ke jam dengan \texttt{/} dan \texttt{\%}, memilih tarif dengan \texttt{if}, dan mengulang untuk beberapa kendaraan.

Hari ini kita berlatih memecah soal seperti itu menjadi bagian-bagian yang sudah kalian kuasai.

\begin{MKeluaran}[title={Tampilan yang diinginkan}]
M 61   -> 2 jam, tarif 3000
B 180  -> 3 jam, tarif 11000
X 10   -> Jenis tidak dikenal
\end{MKeluaran}
\begin{Tebak}[Coba pikirkan]
Materi pertemuan mana saja yang dipakai soal ini? Bagian mana yang menurut kalian paling rawan salah?
\end{Tebak}
\begin{Jawaban}[Cara memecahnya]
Baca soal: Garis bawahi masukan, keluaran, dan aturan. Pecah: Petakan setiap aturan ke konstruksi C++ yang sudah dipelajari. Uji: Siapkan data uji, termasuk nilai batas, sebelum menulis kode.
\end{Jawaban}
\Bagian{Langkah 1}{Peta materi UTS}
Tabel ini merangkum konsep kunci dan jebakan yang paling sering muncul di setiap pertemuan. Gunakan sebagai daftar periksa: untuk setiap baris, pastikan kalian bisa menulis contoh kode dan menjelaskan jebakannya.

\par\vspace{4pt}\noindent{\fontsize{9.5pt}{13pt}\selectfont
\begin{tabular}{@{}>{\raggedright\arraybackslash}p{0.142\textwidth}>{\raggedright\arraybackslash}p{0.392\textwidth}>{\raggedright\arraybackslash}p{0.446\textwidth}@{}}
\kepala{Pertemuan} & \kepala{Konsep kunci} & \kepala{Jebakan yang sering muncul}\\
P1 & struktur program, \texttt{cout}, escape & garis miring tunggal, lupa titik koma\\\garis
P2 & tipe data, \texttt{cin}, \texttt{getline} & pembagian bulat, \texttt{getline} sesudah \texttt{cin}\\\garis
P3 & \texttt{\%}, precedence, \texttt{++}, logika & urutan operasi, \texttt{x++} vs \texttt{++x}\\\garis
P4 & \texttt{if}, \texttt{else if}, \texttt{switch} & \texttt{=} vs \texttt{==}, urutan syarat, \texttt{break}\\\garis
P5 & \texttt{for}, \texttt{while}, \texttt{do-while} & akumulator tidak diinisialisasi, off-by-one\\\garis
P6 & bersarang, \texttt{break}, \texttt{continue} & batas dalam, rumus pola\\
\garisdasar
\end{tabular}}\par\vspace{6pt}
\Bagian{Langkah 2}{Strategi soal A}
Soal A dinilai dari program yang berjalan dan keluaran yang tepat. Strategi yang paling aman adalah membangun program sedikit demi sedikit dan selalu menjaganya dalam keadaan terkompilasi.

\par\vspace{4pt}\noindent{\fontsize{9.5pt}{13pt}\selectfont
\begin{tabular}{@{}>{\raggedright\arraybackslash}p{0.280\textwidth}>{\raggedright\arraybackslash}p{0.700\textwidth}@{}}
\kepala{Bagian} & \kepala{Isi}\\
Rancang dulu & Tulis masukan, proses, keluaran, dan data uji di kertas sebelum mengetik.\\\garis
Bangun bertahap & Baca masukan, tampilkan, kompilasi. Tambah satu bagian, kompilasi lagi.\\\garis
Uji batas & Jalankan dengan data uji biasa dan nilai batas. Cocokkan format keluaran.\\
\garisdasar
\end{tabular}}\par\vspace{6pt}
\begin{Penting}[]
Program yang sebagian benar dan terkompilasi mendapat nilai lebih baik daripada program lengkap yang tidak terkompilasi.
\end{Penting}
\Bagian{Langkah 3}{Strategi soal B}
Soal B dinilai dari ketepatan tracing dan alasan. Trace table yang rapi membuat penilai bisa memberi poin untuk langkah yang benar walaupun ada satu kesalahan di tengah.

\par\vspace{4pt}\noindent{\fontsize{9.5pt}{13pt}\selectfont
\begin{tabular}{@{}>{\raggedright\arraybackslash}p{0.280\textwidth}>{\raggedright\arraybackslash}p{0.700\textwidth}@{}}
\kepala{Bagian} & \kepala{Isi}\\
Buat tabel & Satu kolom per variabel, satu baris per perubahan nilai.\\\garis
Ikuti persis & Kerjakan sesuai aturan C++: pembagian bulat, precedence, \texttt{break}.\\\garis
Jelaskan & Pakai istilah tepat: pembagian bulat, nilai batas, akumulator, cabang.\\
\garisdasar
\end{tabular}}\par\vspace{6pt}
\begin{Penting}[]
Jawaban B tanpa langkah hanya mendapat sebagian poin, walaupun hasil akhirnya benar.
\end{Penting}
\Bagian{Langkah 4}{Contoh tracing gabungan}
Contoh ini menggabungkan perulangan, sisa bagi, dan rantai \texttt{else if}. Kuncinya ada di i = 6: angka ini habis dibagi 2 dan 3, tetapi karena syarat genap diperiksa lebih dulu, cabang kedua tidak pernah dijalankan untuk i = 6.

\begin{MKode}{gabungan.cpp}
int n = 6, hasil = 0;
for (int i = 1; i <= n; i++) {
    if (i % 2 == 0) {
        hasil += i;
    } else if (i % 3 == 0) {
        hasil -= i;
    }
}
cout << hasil << endl;
\end{MKode}
\begin{Tebak}[]
Sebelum menjalankan program, tulis keluarannya di kertas, karakter demi karakter.
\end{Tebak}
\begin{Jawaban}[]
\begin{MKeluaran}[]
9
\end{MKeluaran}
i = 2, 4, 6 menambah. i = 3 mengurangi. i = 6 genap, jadi cabang kedua tidak diperiksa. Hasil: 2 + 4 + 6 − 3 = 9.
\end{Jawaban}
\Bagian{Langkah 5}{Memecah digit}
Pasangan \texttt{\% 10} dan \texttt{/ 10} di dalam \texttt{while} adalah pola klasik untuk memproses digit sebuah bilangan: menjumlahkan, menghitung, atau membalik digit.

\begin{MKode}{digit.cpp}
int n = 4096, jumlah = 0;
while (n > 0) {
    jumlah += n % 10;
    n /= 10;
}
cout << jumlah << endl;
\end{MKode}
\begin{Tebak}[]
Sebelum menjalankan program, tulis keluarannya di kertas, karakter demi karakter.
\end{Tebak}
\begin{Jawaban}[]
\begin{MKeluaran}[]
19
\end{MKeluaran}
\texttt{\% 10} mengambil digit terakhir, \texttt{/ 10} membuangnya. Pola ini muncul di banyak soal.
\end{Jawaban}
\Bagian{Waspada}{Pesan error yang sering muncul}
Semua pesan di tabel ini dihasilkan g++ dengan opsi \texttt{-Wall}, sama seperti yang dipakai \texttt{cek.sh}.
\par\vspace{4pt}\noindent{\fontsize{9.5pt}{13pt}\selectfont
\begin{tabular}{@{}>{\raggedright\arraybackslash}p{0.240\textwidth}>{\raggedright\arraybackslash}p{0.270\textwidth}>{\raggedright\arraybackslash}p{0.470\textwidth}@{}}
\kepala{Kesalahan} & \kepala{Contoh} & \kepala{Pesan pertama dari g++}\\
Lupa titik koma & \texttt{cout << x} & \texttt{error: expected ';' before 'return'}\\\garis
Variabel belum dideklarasikan & \texttt{total += x;} & \texttt{error: 'total' was not declared in this scope}\\\garis
= di dalam kondisi & \texttt{if (a = b)} & \texttt{warning: suggest parentheses around assignment used as truth value}\\\garis
break di luar perulangan & \texttt{break; di if} & \texttt{error: break statement not within loop or switch}\\\garis
Kurung kurawal kurang & \texttt{while \{ if \{ \}} & \texttt{error: expected '\}' at end of input}\\
\garisdasar
\end{tabular}}\par\vspace{6pt}
Kelima error ini mewakili kesalahan yang paling sering membuat program UTS tidak terkompilasi. Biasakan mengompilasi setiap beberapa baris agar error cepat ditemukan.

\Bagian{Latihan}{Latihan mandiri}
Latihan ini sama dengan hands-on di kelas. Semua file ada di \texttt{02 Hands-on/P7 - Review UTS - Hands-on/latihan/}. Kerjakan berurutan, lalu cocokkan dengan \texttt{expected/} atau jalankan \texttt{bash cek.sh}.
\par\vspace{4pt}\noindent{\fontsize{9.5pt}{13pt}\selectfont
\begin{tabular}{@{}>{\raggedright\arraybackslash}p{0.070\textwidth}>{\raggedright\arraybackslash}p{0.360\textwidth}>{\raggedright\arraybackslash}p{0.150\textwidth}>{\raggedright\arraybackslash}p{0.400\textwidth}@{}}
\kepala{No} & \kepala{File} & \kepala{Tingkat} & \kepala{Yang dilatih}\\
1 & \texttt{handson1-parkir.cpp} & Dasar & masukan, \texttt{/} dan \texttt{\%}, \texttt{else if}\\\garis
2 & \texttt{handson2-digit.cpp} & Menengah & \texttt{while} dengan pola digit\\\garis
3 & \texttt{handson3-telusur.cpp} & Tantangan & trace table bersarang dengan \texttt{continue} dan \texttt{break}\\\garis
+ & \texttt{bonus-menu.cpp} & Pengayaan & menu \texttt{do-while}, validasi\\
\garisdasar
\end{tabular}}\par\vspace{6pt}
\Sub{Latihan 1: handson1-parkir.cpp}
Tarif parkir motor dan mobil dengan pembulatan jam ke atas. Kerjakan tanpa bantuan sebagai simulasi soal A.

\Sub{Latihan 2: handson2-digit.cpp}
Hitung banyak digit, jumlah digit, digit terbesar, dan bilangan dibalik.

\Sub{Latihan 3: handson3-telusur.cpp}
Simulasi soal B: tulis trace table dan perkiraan keluaran sebelum menjalankan. File ini sudah lulus \texttt{cek.sh}; yang dinilai adalah tracing kalian.

\Sub{Bonus: bonus-menu.cpp}
ATM mini dengan cek saldo, setor, dan tarik yang divalidasi.

\Sub{Latihan menelusuri}
\begin{MKode}{tebak.cpp}
int a = 5, b = 0;
while (a > 0) {
    if (a % 2 == 1) {
        b += a;
    } else {
        b -= 1;
    }
    a -= 2;
}
cout << a << " " << b << endl;
\end{MKode}
\begin{Tebak}[]
Tulis keluarannya di kertas tanpa menjalankan program. Buat trace table bila perlu.
\end{Tebak}
\begin{Jawaban}[]
\begin{MKeluaran}[]
-1 9
\end{MKeluaran}
\texttt{a} bernilai 5, 3, 1, semuanya ganjil, jadi \texttt{b} = 5 + 3 + 1 = 9. Sesudah itu \texttt{a} = −1 dan perulangan berhenti.
\end{Jawaban}
\Bagian{Penutup}{Rangkuman}
\par\vspace{4pt}\noindent{\fontsize{9.5pt}{13pt}\selectfont
\begin{tabular}{@{}>{\raggedright\arraybackslash}p{0.320\textwidth}>{\raggedright\arraybackslash}p{0.660\textwidth}@{}}
\kepala{Konsep} & \kepala{Yang perlu diingat}\\
Peta materi UTS & Tabel ini merangkum konsep kunci dan jebakan yang paling sering muncul di setiap pertemuan.\\\garis
Strategi soal A & Program yang sebagian benar dan terkompilasi mendapat nilai lebih baik daripada program lengkap yang tidak terkompilasi.\\\garis
Strategi soal B & Jawaban B tanpa langkah hanya mendapat sebagian poin, walaupun hasil akhirnya benar.\\\garis
Contoh tracing gabungan & i = 2, 4, 6 menambah.\\\garis
Memecah digit & \texttt{\% 10} mengambil digit terakhir, \texttt{/ 10} membuangnya.\\
\garisdasar
\end{tabular}}\par\vspace{6pt}
\Bagian{Penutup}{Cek pemahaman}
Jawab tanpa menjalankan program, lalu periksa dengan menjalankannya.
\begin{enumerate}
\item Berapa jam untuk 120 menit dan 121 menit dengan pembulatan ke atas?
\item Apa keluaran \texttt{int n = 305; cout << n \% 10 + n / 100;}?
\item Pada rantai \texttt{else if}, apakah dua cabang bisa jalan bersamaan?
\item Berapa kali \texttt{cout} dijalankan oleh \texttt{for i 1..3} berisi \texttt{for j 1..i}?
\item Sebutkan tiga isi minimal sebuah trace table.
\end{enumerate}
\begin{Jawaban}[]
\begin{enumerate}[leftmargin=16pt]
\item 2 dan 3 jam.
\item 8 (5 + 3).
\item Tidak, hanya cabang pertama yang syaratnya benar.
\item 6 kali (1 + 2 + 3).
\item Kolom variabel, baris setiap perubahan, dan keluaran.
\end{enumerate}
\end{Jawaban}
\Bagian{Penilaian}{Exit Ticket 7}
Dikerjakan di 25 menit terakhir kelas, sendiri, tanpa AI, internet, atau diskusi. Exit Ticket 7 adalah satu dari 14 Exit Ticket; rata-ratanya menjadi komponen berbobot 25\%.
\par\vspace{4pt}\noindent{\fontsize{9.5pt}{13pt}\selectfont
\begin{tabular}{@{}>{\raggedright\arraybackslash}p{0.080\textwidth}>{\raggedright\arraybackslash}p{0.600\textwidth}>{\raggedright\arraybackslash}p{0.100\textwidth}>{\raggedright\arraybackslash}p{0.200\textwidth}@{}}
\kepala{Soal} & \kepala{Isi} & \kepala{Poin} & \kepala{CPMK}\\
A1 & Tulis program singkat yang memakai masukan, percabangan, dan perulangan & 25 & CPMK0607\\\garis
A2 & Perbaiki program agar keluarannya sesuai contoh & 25 & CPMK0607\\\garis
B1 & Buat trace table untuk perulangan bersarang dengan \texttt{continue} & 20 & CPMK0608\\\garis
B2 & Jelaskan mengapa dua cabang tidak dijalankan bersamaan pada sebuah rantai \texttt{else if} & 20 & CPMK0608\\\garis
B3 & Sebutkan materi pertemuan yang dipakai sebuah soal dan alasannya & 10 & CPMK0608\\
\garisdasar
\end{tabular}}\par\vspace{6pt}
\Bagian{Penutup}{Bacaan lanjut dan persiapan minggu depan}
\begin{enumerate}
\item Modul Belajar 1 sampai 6.
\item Deitel, P. dan Deitel, H. C++ How to Program, edisi ke-10, Bab 2 sampai 5.
\item Slide \texttt{01 Slide/P7 Review dan Simulasi UTS.pdf} dan hands-on di \texttt{02 Hands-on/P7 - Review UTS - Hands-on/}.
\end{enumerate}
\begin{Penting}[Minggu depan]
Pertemuan 8 adalah UTS. Kisi-kisi, tata tertib, dan contoh soal ada di slide P8 dan Modul Belajar 8.
\end{Penting}
\end{document}
